%******************************************************* Abstract
%*******************************************************
%\renewcommand{\abstractname}{Abstract}
\pdfbookmark[1]{Abstract}{Abstract} \begingroup \let\clearpage\relax
\let\cleardoublepage\relax \let\cleardoublepage\relax

\chapter*{Abstract}

Privacy-preserving systems are often designed as if the choice is between
decentralised privacy and centralised surveillance. In practice, systems that
protect users already rely on issuers, auditors, service providers and central
banks, which provide availability and accountability and are often constrained
by legal frameworks. This thesis argues that practical privacy comes from
modelling these participants and their responsibilities, and using cryptography
to limit what each one can learn and do. An institution may
perform computation, maintain state, issue credentials, validate transactions
or provide availability without being trusted with users' private information.

We develop four protocols, each starting from a role that exists in practice.
First, we give a framework for threshold structure-preserving signature
schemes, which let several parties jointly sign commitments, with all
interaction between signers moved to an offline phase. Second, we construct
multi-issuer aggregate anonymous credentials, in which credentials from
independent issuers are combined and presented in a single proof. A
registration authority binds credentials to the same user without learning
their attributes. Third, we give a transaction tracing protocol in
which an auditor's capability is granted explicitly and scoped to one user and
one period of time, and the auditor keeps no persistent per-user state. Fourth,
we construct a transferable offline payment scheme for central bank digital currencies. A
secure element prevents double spending, double spenders are identified if the
secure element is broken, and payments by honest users remain unlinkable.

Each protocol comes with a generic construction from standard primitives, an
efficient instantiation, a security proof and an implementation. Auditing a period
of a million transactions takes 37 ms, and a smart card acting as the secure
element authorises an offline payment in under 400 ms. Our protocols balance
users' privacy against the requirements of existing infrastructure, such as
regulation, accountability and availability. Adding auditing to a
privacy-preserving protocol can seem to defeat its purpose, but the alternative
to auditable payments is not perfect privacy. It is the status quo, where banks
and payment providers see all transactions in the clear.

\endgroup

\cleardoublepage%

\begingroup \let\clearpage\relax \let\cleardoublepage\relax
\let\cleardoublepage\relax

\pdfbookmark[1]{Resumé}{Résumé} \chapter*{Résumé}

\begin{otherlanguage}{french}
Les systèmes respectueux de la vie privée sont souvent conçus comme s'il
fallait choisir entre une confidentialité décentralisée et une surveillance
centralisée. En pratique, les systèmes qui protègent leurs utilisateurs reposent
déjà sur des émetteurs, des auditeurs, des fournisseurs de services et des
banques centrales, qui assurent la disponibilité et la responsabilité et sont
souvent soumis à des cadres juridiques. Cette thèse soutient qu'une
confidentialité praticable vient d'une modélisation de ces participants et de
leurs responsabilités, et de l'usage de la cryptographie pour limiter ce que
chacun d'eux peut apprendre et faire. Une institution peut effectuer des
calculs, maintenir un état, émettre des attestations, valider des transactions
ou assurer la disponibilité sans qu'on lui confie les informations privées des
utilisateurs.

Nous développons quatre protocoles, chacun partant d'un rôle qui existe en
pratique. Premièrement, nous proposons un cadre pour les schémas de signature à
seuil préservant la structure, qui permettent à plusieurs parties de signer
conjointement des engagements, toute interaction entre les signataires ayant
lieu dans une phase hors ligne. Deuxièmement, nous construisons des
attestations anonymes agrégées à émetteurs multiples, dans lesquelles les
attestations d'émetteurs indépendants sont combinées et présentées en une seule
preuve. Une autorité d'enregistrement lie les attestations à un même
utilisateur sans apprendre ses attributs. Troisièmement, nous proposons un
protocole de traçage des transactions dans lequel la capacité d'un
auditeur est accordée explicitement et limitée à un utilisateur et à une
période, et l'auditeur ne conserve aucun état persistant par utilisateur. Quatrièmement, nous
construisons un schéma de paiement hors ligne transférable pour les monnaies
numériques de banque centrale. Un élément sécurisé empêche la double dépense,
les auteurs d'une double dépense sont identifiés si l'élément sécurisé est
compromis, et les paiements des utilisateurs honnêtes restent non reliables
entre eux.

Chaque protocole s'accompagne d'une construction générique à partir de
primitives standard, d'une instanciation efficace, d'une preuve de sécurité et
d'une implémentation. L'audit d'une période d'un million de
transactions prend 37~ms, et une carte à puce jouant le rôle d'élément sécurisé
autorise un paiement hors ligne en moins de 400~ms. Nos protocoles concilient
la vie privée des utilisateurs avec les exigences des infrastructures
existantes, comme la réglementation, la responsabilité et la disponibilité.
Ajouter un audit à un protocole respectueux de la vie privée peut sembler
contraire à son objectif, mais l'alternative aux paiements auditables n'est pas
une confidentialité parfaite. C'est le statu quo, où les banques et les
prestataires de paiement voient toutes les transactions en clair.
\end{otherlanguage}

\endgroup

\vfill
